\section*{Capítulo \ref{ch:b3:cancer-biology} --- Biología del cáncer}

\begin{solution}{exo:b3:cancer-biology:1}
Un \omterm{def:b3:cancer-biology:genes}{oncogén} es una versión con ganancia de función de un gen que favorece
la proliferación o la supervivencia; basta un alelo mutante (es
dominante): \emph{RAS} (mutación puntual), \emph{MYC} (translocación,
amplificación), \emph{HER2}, \emph{BCR--ABL}. Un supresor de \omterm{def:b3:cancer-biology:hallmarks}{tumores}
frena la proliferación o impone la parada, la muerte o la reparación; han
de perderse los dos alelos (es recesivo a nivel celular): \emph{RB},
\emph{TP53}, \emph{APC}, \emph{BRCA1}.
\end{solution}

\begin{solution}{exo:b3:cancer-biology:2}
Proliferación sostenida: \emph{RAS}, \emph{MYC}. Insensibilidad a los
supresores del crecimiento: \emph{RB}, \emph{CDKN2A} (p16). Resistencia a
la muerte: \emph{BCL2}, \emph{TP53}. Inmortalidad replicativa: el
promotor de la \omterm{def:b3:cancer-biology:telomerase}{telomerasa}. Angiogénesis: \emph{VEGF} inducido a través de
HIF (pérdida de \emph{VHL}). \omterm{def:b3:cancer-biology:metastasis}{Invasión} y \omterm{def:b3:cancer-biology:hallmarks}{metástasis}: pérdida de
E-cadherina, los factores de la \omterm{def:b3:cancer-biology:metastasis}{transición epitelio--mesénquima}.
Inestabilidad del genoma: los genes de la \omterm{prop:b3:dna-repair:fidelity}{reparación de errores de
apareamiento}, \emph{\omterm{def:b3:dna-repair:dsb}{BRCA}}. Metabolismo desregulado: \emph{MYC}, la \omterm{prop:b3:cancer-biology:pathways}{vía
PI3K}. \omterm{prop:b3:cancer-biology:immune}{Evasión inmunitaria}: \emph{PD-L1}, pérdida del MHC. Inflamación
favorecedora del \omterm{def:b3:cancer-biology:hallmarks}{tumor}: la señalización de \emph{NF-$\kappa$B}.
\end{solution}

\begin{solution}{exo:b3:cancer-biology:3}
Los niños con antecedentes familiares desarrollaban varios \omterm{def:b3:cancer-biology:hallmarks}{tumores} en los
dos ojos, y pronto; los casos esporádicos desarrollaban un solo \omterm{def:b3:cancer-biology:hallmarks}{tumor}, y
más tarde. Infirió que hacían falta dos sucesos en una célula: los
pacientes familiares habían heredado uno y solo necesitaban un segundo
suceso somático (lo bastante frecuente como para darse en varias células
retinianas), mientras que los esporádicos necesitaban que los dos
ocurrieran de forma somática en una misma célula --- raro, tardío, único.
De ahí un gen cuyas dos copias han de perderse ambas: el primer supresor
de \omterm{def:b3:cancer-biology:hallmarks}{tumores}.
\end{solution}

\begin{solution}{exo:b3:cancer-biology:4}
Se siembra sin histidina, junto con el compuesto de prueba, un mutante de
\emph{Salmonella} que necesita histidina; solo surgen colonias de células
en las que una mutación nueva ha revertido el defecto, de modo que su
número mide la mutagenicidad. Se añade extracto de hígado porque muchos
carcinógenos son inocuos hasta que las enzimas hepáticas los convierten
en metabolitos reactivos, como ocurre en el cuerpo; las bacterias carecen
de esas enzimas.
\end{solution}

\begin{solution}{exo:b3:cancer-biology:5}
La incidencia sube $100$ veces mientras la edad se duplica:
$2^{k-1} = 100$, $k - 1 =
\log_{2}100 = 6.6$, luego unos siete u ocho sucesos limitantes en la
lectura ingenua --- menos si se admite la expansión clonal entre golpes.
\end{solution}

\begin{solution}{exo:b3:cancer-biology:6}
$L = \sqrt{2Dc_{0}/q}$: para $q = 0.01$, $\sqrt{2\times 2\times 10^{-9}
\times 0.05/0.01} = \qty{141}{\micro m}$; para $q = 0.03$,
\qty{82}{\micro m}. Un \omterm{def:b3:cancer-biology:hallmarks}{tumor} de crecimiento rápido respira más deprisa,
agota el oxígeno en una distancia menor y su centro muere a un tamaño más
pequeño.
\end{solution}

\begin{solution}{exo:b3:cancer-biology:7}
$10^{11}\times 10^{-3n} < 1$ exige $n > 3.7$: cuatro tandas. Tras dos
tandas quedan $10^{5}$ células --- una décima de milímetro cúbico, muy
por debajo de las $10^{9}$ que pueden detectarse ---, de modo que el
\omterm{def:b3:cancer-biology:hallmarks}{tumor} ha ``desaparecido'' para toda prueba mientras quedan cien mil
células; el tratamiento ha de continuar hasta el número de tandas
previsto, no hasta que el \omterm{def:b3:cancer-biology:hallmarks}{tumor} desaparezca.
\end{solution}

\begin{solution}{exo:b3:cancer-biology:8}
Las dos lesiones activan la misma vía (receptor $\to$ \omterm{def:b3:cancer-biology:genes}{Ras} $\to$ ERK); una
vez presente una, la otra no confiere ninguna ventaja adicional y no se
selecciona. Un inhibidor de EGFR bloquea el receptor; una célula que
adquiere una mutación de KRAS aguas abajo restaura la señal con el
receptor todavía bloqueado, y queda seleccionada durante el tratamiento
--- resistencia por desvío.
\end{solution}

\begin{solution}{exo:b3:cancer-biology:9}
E7 libera a E2F de Rb y lleva a la célula epitelial infectada más allá
del \omterm{def:b3:cell-cycle-apoptosis:checkpoints}{punto de restricción} sin factores de crecimiento; E6 destruye a \omterm{def:b3:cancer-biology:genes}{p53},
de modo que la proliferación inadecuada y el daño que causa no
desencadenan ni parada ni \omterm{def:b3:cell-cycle-apoptosis:apoptosis}{apoptosis}. Se aportan dos rasgos distintivos de
una vez, pero los restantes --- la inmortalidad, la \omterm{def:b3:cancer-biology:metastasis}{invasión}, la
inestabilidad --- necesitan más mutaciones somáticas, que tardan décadas
en acumularse en las células infectadas de forma persistente. La vacuna
provoca anticuerpos neutralizantes contra la cápside y previene la
infección; administrada antes de la exposición, nunca hay una célula
infectada que transformar.
\end{solution}

\begin{solution}{exo:b3:cancer-biology:10}
Tras el $i$-ésimo golpe el clon crece por $g$, de modo que hay $g^{i}$
veces más células en riesgo del golpe siguiente; la probabilidad de la
secuencia completa en el linaje descendiente de una célula original es
$g\cdot g^{2}\cdots$ --- para un $g$ constante por golpe el riesgo
acumulado pasa a ser $N(ut)^{k}
g^{k-1}$ salvo factores que dependen del orden, y la incidencia conserva
la potencia $t^{k-1}$ con un prefactor $g^{k-1}$ mayor. Si las propias
expansiones crecen con el tiempo (clones de crecimiento exponencial), la
curva se empina aún más, de modo que una pendiente de $6$ la producen
menos de siete conductoras: el $k$ de Armitage--Doll es una cota superior
del número de conductoras limitantes, y los \omterm{def:b3:cancer-biology:hallmarks}{tumores} secuenciados muestran
de dos a ocho.
\end{solution}

\begin{solution}{exo:b3:cancer-biology:11}
(a) El fármaco A solo se enfrenta a $\mu_{A}N = 100$ células resistentes
preexistentes: $P_{0} = e^{-100} \approx 0$; el clon resistente a A
sobrevive y vuelve a crecer durante las dieciocho semanas de A, y cuando
empieza B es de nuevo una población grande, de modo que B se enfrenta
también a $\mu_{B}N' \gg 1$ --- la secuencia fracasa dos veces.
(b) Juntos: una célula doblemente resistente surge a $10^{-14}$ por
división, $\mu_{A}\mu_{B}N = 10^{-5}$, $P_{0} =
0.99999$. La combinación desde el principio, cuando $N$ es menor, es la
norma porque lo que hace improbable la resistencia es el producto de dos
velocidades pequeñas.
\end{solution}

\begin{solution}{exo:b3:cancer-biology:12}
Con $100$ veces más células y el mismo $u$ y el mismo $k$, el riesgo
acumulado $N(ut)^{k}$ sería $100$ veces mayor --- todos los elefantes
deberían morir de \omterm{def:b3:cancer-biology:hallmarks}{cáncer} siendo jóvenes. No ocurre, luego $u$ o $k$ han
de ser distintos: veinte copias de \emph{TP53} hacen más fuerte la
respuesta apoptótica al daño y retiran las células mutantes antes de que
se expandan (lo que baja la $u$ efectiva de ese paso); los animales
grandes tienen además tasas de mutación por célula más bajas (mejor
reparación, menos divisiones por célula y año porque sus células se
recambian más despacio) y pueden exigir más golpes (una capa supresora
adicional). La supresión del \omterm{def:b3:cancer-biology:hallmarks}{cáncer} evolucionó junto con el tamaño
corporal.
\end{solution}

\begin{solution}{pb:b3:cancer-biology:1}
\textbf{1.} $10^{12}/10^{3} = 10^{9}$ células; $\log_{2}10^{9} \approx 30$
duplicaciones.
\textbf{2.} $30\times 100 = 3000$ días, unos $8$ años.
\textbf{3.} \qty{1}{kg} $\approx 10^{12}$ células: diez duplicaciones
más, $1000$ días, menos de tres años --- un tercio del tiempo hasta la
detección.
\textbf{4.} El tiempo de duplicación inicial debió de ser de unos $60$
días, o bien el \omterm{def:b3:cancer-biology:hallmarks}{tumor} creció más deprisa cuando era pequeño y se frenó al
crecer (el patrón habitual, a medida que fallan el suministro y el
espacio).
\textbf{5.} Un ciclo de $2$ días duplicaría la población cada $2$ días;
un tiempo de duplicación de $100$ significa que el crecimiento neto es el
$2\,\%$ de la velocidad de nacimiento: casi cada célula que nace la
compensa una que muere, o bien pocas están ciclando --- nacimiento y
muerte casi equilibrados, con el crecimiento como diferencia pequeña.
\textbf{6.} Los fármacos matan células que ciclan, de modo que un \omterm{def:b3:cancer-biology:hallmarks}{tumor}
con poca fracción en ciclo pierde una fracción menor por tanda; pero el
\omterm{def:b3:cancer-biology:hallmarks}{tumor} rápido, que ha perdido más, también vuelve a crecer entre tandas.
En balance, los \omterm{def:b3:cancer-biology:hallmarks}{tumores} rápidos (leucemias, linfomas, \omterm{def:b3:cancer-biology:hallmarks}{cáncer} de
testículo) son los que la quimioterapia cura.
\textbf{7.} Razón $300$ sobre una razón de edades $80/30 = 2.67$: $k - 1 =
\ln 300/\ln 2.67 = 5.8$, $k \approx 7$.
\textbf{8.} $N(ut)^{k} = 10^{8}\times (7\times 10^{-5})^{7} \approx
10^{-21}$: absurdo frente a un riesgo de por vida cercano a $1$ de cada
$3$. Falta: la expansión clonal que multiplica las células en riesgo tras
cada golpe, las tasas de mutación elevadas por la inestabilidad, muchas
más divisiones (células madre ciclando durante décadas) y menos
conductoras verdaderas.
\textbf{9.} Multiplicar por $g^{k-1} = 100^{6} = 10^{12}$ da $10^{-9}$
--- todavía minúsculo con siete sucesos; con cuatro sucesos y
expansiones, $10^{8}\times (7\times 10^{-5})^{4}\times 10^{6} \approx 2\times
10^{-3}$, el orden correcto. El cuadro real es de pocas conductoras y
grandes expansiones.
\textbf{10.} $(ut)^{k-1}/(ut)^{k} = 1/(ut) = 1/(10^{-6}\times 40) =
2.5\times 10^{4}$: la incidencia de una portadora a los cuarenta es
decenas de miles de veces la esporádica --- los \omterm{def:b3:cancer-biology:hallmarks}{cánceres} hereditarios
golpean jóvenes.
\textbf{11.} Un suceso diez veces más rápido: incidencia $\times 10$;
dos: $\times 100$. El riesgo observado, de quince a veinticinco veces,
sugiere que el humo actúa sobre más de un paso.
\textbf{12.} Al dejarlo, $u$ vuelve a bajar para los sucesos que quedan
por venir, de modo que la incidencia deja de subir tan empinadamente;
pero los golpes ya acumulados en las células siguen ahí, de modo que las
células del exfumador están más cerca del \omterm{def:b3:cancer-biology:hallmarks}{cáncer} que las de quien nunca
fumó y el exceso de riesgo persiste, congelado más que borrado.
\textbf{13.} $L = \sqrt{2\times 2\times 10^{-9}\times 0.05/0.03} =
\qty{82}{\micro m}$.
\textbf{14.} Totalmente oxigenado hasta un diámetro $2L \approx
\qty{160}{\micro m}$. Un esferoide de \qty{2}{mm} tiene un borde vivo de
\qty{82}{\micro m}: fracción viva $1 - (918/1000)^{3} = 0.23$.
\textbf{15.} Las células de en medio están a \qty{75}{\micro m} de un
vaso, justo dentro de $L$ --- oxigenadas en principio, pero los vasos
tumorales llevan menos oxígeno y fluyen de forma intermitente, de modo
que esas células son hipóxicas. Las células hipóxicas son de dos a tres
veces más resistentes a la radiación (hace falta oxígeno para que el daño
se vuelva permanente), y la hipoxia desencadena \omterm{def:b3:cell-cycle-apoptosis:apoptosis}{apoptosis} dependiente de
\omterm{def:b3:cancer-biology:genes}{p53}, de modo que selecciona a las células que han perdido \omterm{def:b3:cancer-biology:genes}{p53}.
\textbf{16.} Reducir la densidad a la mitad separa los vasos en
$\sqrt{2}$: quedan a \qty{212}{\micro m}, con las células de en medio a
\qty{106}{\micro m}, más allá de $L$: mueren o, si sobreviven hipóxicas,
se vuelven más agresivas y angiogénicas.
\textbf{17.} $30/2 = 15$ veces más glucosa. Los \omterm{def:b3:cancer-biology:hallmarks}{tumores} concentran por
tanto los análogos de la glucosa, y una tomografía por emisión de
positrones con fluorodesoxiglucosa los ilumina.
\textbf{18.} Los \omterm{def:b3:cancer-biology:hallmarks}{tumores} hambrientos se encogen hasta el tamaño limitado
por la difusión y persisten, latentes e hipóxicos, o se apropian de vasos
existentes; la hipoxia selecciona clones agresivos. Pero podar la
vasculatura caótica del \omterm{def:b3:cancer-biology:hallmarks}{tumor} baja la presión intersticial y uniformiza
el flujo, de modo que la quimioterapia penetra mejor --- la combinación
funciona donde ninguna de las dos lo hace por separado.
\textbf{19.} $10^{9}\times 10^{-2n} < 1$: $n > 4.5$, cinco tandas.
\textbf{20.} $21$ días $= 0.21$ duplicaciones, recrecimiento
$\times 1.16$; factor neto por ciclo $0.0116$;
$\ln 10^{-9}/\ln 0.0116 = 4.65$: siguen siendo cinco ciclos.
\textbf{21.} $10^{-7}\times 10^{9} = 100$ células resistentes esperadas;
$P_{0} = e^{-100} \approx 0$: la resistencia es segura.
\textbf{22.} $\mu_{1}\mu_{2}N = 10^{-14}\times 10^{9} = 10^{-5}$;
$P_{0} = 0.99999$. Con $10^{12}$ células, $10^{-2}$, $P_{0} = 0.99$.
\textbf{23.} Con $10^{6}$ células: un fármaco, $\mu N = 0.1$,
$P_{0} = 0.90$; dos fármacos, $10^{-8}$, $P_{0} \approx 1$. Trátese
cuando la población es pequeña, y trátese con combinaciones.
\textbf{24.} Los linfocitos~T atacan muchos neoantígenos a la vez, de
modo que ninguna mutación aislada se les escapa; escapar exige perder la
presentación de antígenos misma (MHC, $\beta_{2}$-microglobulina) o
suprimir a los linfocitos~T, una clase de suceso distinta y más rara. La
misma carga alta de mutaciones que hace segura la resistencia de un \omterm{def:b3:cancer-biology:hallmarks}{tumor}
a un inhibidor de quinasas le da muchos neoantígenos: los \omterm{def:b3:cancer-biology:hallmarks}{tumores}
deficientes en \omterm{prop:b3:dna-repair:fidelity}{reparación de errores de apareamiento} y los inducidos por
el ultravioleta y por el humo son los que mejor responden.
\textbf{25.} Unos $8$ años hasta la detección; \omterm{def:b3:cancer-biology:hallmarks}{tumor} vivo hasta
\qty{82}{\micro m} alrededor de un vaso; y $P_{0} = 0.99999$ de que una
combinación de dos fármacos no encuentre ninguna célula resistente en un
\omterm{def:b3:cancer-biology:hallmarks}{tumor} de \qty{1}{cm^3}.
\end{solution}
